\documentclass[11pt,a4paper]{article}
\usepackage[margin=23mm]{geometry}
\usepackage[T1]{fontenc}
\usepackage[utf8]{inputenc}
\usepackage{lmodern,microtype}
\usepackage{amsmath,amssymb,booktabs,longtable,array,tabularx,xcolor,fancyhdr}
\usepackage[hidelinks]{hyperref}
\numberwithin{equation}{section}
\pagestyle{fancy}\fancyhf{}
\lhead{Long-Only Capacity Methodology}
\rhead{Revised 27 September 2026}
\cfoot{\thepage}\setlength{\headheight}{14pt}
\newcolumntype{Y}{>{\raggedright\arraybackslash}X}
\title{\textbf{Long-Only Equity Capacity Methodology}\\[0.4em]
\large Distributional Liquidity Migration, Execution Horizons and Alpha Decay}
\author{Rodrigo Dupleich\thanks{with the help of ChatGTP}}
\date{27 September 2026}

\begin{document}
\maketitle
\begin{abstract}
This document defines the stable methodology for estimating the capacity of a long-only equity
strategy. It separates the methodology from portfolio-specific runs. Each run should produce a
separate report containing the supplied inputs, fitted parameters, diagnostics, scenarios,
thresholds and warnings. Appendices A--C define the user input contract and distinguish supplied,
policy and fitted parameters.
\end{abstract}
\tableofcontents
\clearpage

\section{Framework overview}
The framework maps assets under management (AUM) into a distribution of parent-order
participation, execution horizons, implementation costs and captured alpha:
\begin{equation}\label{eq:flow}
 A\rightarrow F_A(p)\rightarrow D(p)\rightarrow
 \{\bar c(A),R_g(A)\}\rightarrow R_{\mathrm{net}}(A),IR_{\mathrm{net}}(A).
\end{equation}
Turnover $T$ is one-way annual turnover, so executed purchases and sales are approximated by
$2TA$. Cost is expressed in basis points per executed notional.

The model consumes either aggregate ADV buckets or, preferably, parent-order data. Aggregate
bucket shares are interpreted as shares of executed notional, not trade-count frequencies.

\section{Strategy economics}
At current AUM $A_0$, the key quantities are gross relative return $R_g$, tracking error $TE$,
turnover $T$ and average holdings $N$:
\begin{equation}\label{eq:economics}
 IR_g=\frac{R_g}{TE},\qquad
 \alpha_{\mathrm{holding}}=\frac{R_g}{N},\qquad
 \alpha_{\mathrm{cycle}}=\frac{R_g}{T},\qquad
 H\approx\frac{1}{T}.
\end{equation}
Current gross alpha is treated as already reflecting the current implementation schedule. Alpha
decay at alternative AUM levels is therefore measured relative to the current schedule.

\section{Distributional liquidity migration}\label{sec:migration}
Capacity depends on parent-order participation $p=Q/\mathrm{ADV}$, not ADV in isolation. Modelling
ADV alone would require a joint model for parent-order size and its dependence on ADV.

The aggregate implementation fits a Burr XII distribution:
\begin{equation}\label{eq:burr}
 F(p;c,d,\lambda)=1-\left[1+\left(\frac{p}{\lambda}\right)^c\right]^{-d}.
\end{equation}
It offers a flexible body and Pareto-like survival tail $1-F(p)\sim p^{-cd}$. Empirical work finds
power-law tails in trade-size distributions \cite{rak2007}. The calibration constrains $cd\geq
1.25$, preserving a heavy tail while ensuring finite mean participation.

The distribution changes with AUM through
\begin{equation}\label{eq:scaling}
 \lambda(A)=\lambda_0\left(\frac{A}{A_0}\right)^\eta,
 \qquad
 d(A)=d_0\left(\frac{A}{A_0}\right)^{-\kappa}.
\end{equation}
The scale elasticity $\eta$ captures how participation grows with AUM. Tail drift $\kappa$ is
reserved for structural changes in liquidity composition. Bucket shares follow smoothly from
\begin{equation}\label{eq:weights}
 w_i(A)=F_A(u_i)-F_A(\ell_i).
\end{equation}


The participation distribution enters the capacity calculation through its quantile representation.
The distribution is weighted by executed notional: each equal-width quantile interval represents an equal fraction of
traded value, not an equal number of trades. At each AUM, set
$F_A(p)=F(p;c,d(A),\lambda(A))$. Its inverse turns a fixed rank $q\in(0,1)$ into participation:
\begin{equation}\label{eq:quantile}
 p_A(q)=F_A^{-1}(q)
 =\lambda(A)\left[(1-q)^{-1/d(A)}-1\right]^{1/c}.
\end{equation}
The same rank is used at current and alternative AUM. This is a rank-preserving scenario
assumption, not an observation that the same individual orders recur. It supplies the paired
participations needed for the horizon ratio in equation~\eqref{eq:effective}. When $\kappa=0$,
$p_A(q)=p_{A_0}(q)(A/A_0)^\eta$; with tail drift, the full inverse in
equation~\eqref{eq:quantile} is required.

Bucket weights in equation~\eqref{eq:weights} summarize this distribution for charts. The
continuous model does not multiply them into its costs a second time: the notional weights
are already represented by the quantile average made explicit in Section~\ref{sec:net}.

\section{Impact model}
Expected and realised costs are fitted separately using
\begin{equation}\label{eq:impact}
 c(p)=a+bp^\gamma,
\end{equation}
subject to positive parameter bounds. Square-root impact is an important benchmark, while the
implementation permits the exponent to be fitted \cite{farmer2011}. At current AUM, the continuous
distributional expectation is calibrated back to the observed notional-weighted bucket cost.

\section{Execution horizon}
For target daily participation $\rho$ and maximum horizon $D_{\max}$,
\begin{equation}\label{eq:days}
 D(p)=\min\!\left(D_{\max},\max\!\left(1,
 \left\lceil\frac{p}{\rho}\right\rceil\right)\right).
\end{equation}
Current costs already embed the current schedule, so counterfactual effective participation is
\begin{equation}\label{eq:effective}
 p_{\mathrm{eff},j}(A)=p_j(A)\frac{D_j(A_0)}{D_j(A)}.
\end{equation}
The rule exposes the cost-versus-delay trade-off studied in optimal execution
\cite{almgren2001}. A production version can replace it with direct optimisation.

\section{Alpha decay}
Let $H_\alpha$ be signal half-life in trading days and $\delta=\log(2)/H_\alpha$. For equal daily
slices over $D$ days, average alpha capture is
\begin{equation}\label{eq:capture}
 g(D,H_\alpha)=\frac{1}{D}\sum_{t=0}^{D-1}e^{-\delta t}
 =\frac{1-e^{-\delta D}}{D(1-e^{-\delta})}.
\end{equation}
Normalising to current implementation gives
\begin{equation}\label{eq:gross}
 R_g(A)=R_g(A_0)\,
 \mathbb{E}\left[\frac{g(D_j(A),H_{\alpha,j})}
 {g(D_j(A_0),H_{\alpha,j})}\right].
\end{equation}
Aggregate runs currently weight alpha by executed notional. Trade-level runs should weight by each
order's forecast alpha contribution. Predictors with different mean-reversion speeds imply
different desired trading speeds \cite{garleanu2013}.

\section{Net capacity}\label{sec:net}
The participation distribution specified in Section~\ref{sec:migration} determines $p_A(q)$
through equations~\eqref{eq:burr}, \eqref{eq:scaling} and \eqref{eq:quantile}.
For each notional rank $q$, the execution rule in equation~\eqref{eq:days} is evaluated
at current and alternative AUM:
\begin{equation}\label{eq:paired}
 \begin{aligned}
 D_A(q)&=D\!\left(p_A(q)\right),\qquad D_0(q)=D\!\left(p_{A_0}(q)\right),\\
 p_{\mathrm{eff},A}(q)&=p_A(q)\frac{D_0(q)}{D_A(q)}.
 \end{aligned}
\end{equation}
Thus distributional migration has two channels: it changes effective participation and
therefore impact, and it changes execution days and therefore alpha capture.

Let $\bar c_{\mathrm{obs},0}=\sum_i w_{i,0}^{\mathrm{obs}}c_i^{\mathrm{obs}}$ be the observed
current notional-weighted cost in basis points. The impact function $c(p)$ in
equation~\eqref{eq:impact} is the raw fitted curve. Its baseline calibration multiplier is
\begin{equation}\label{eq:calibration}
 K=\frac{\bar c_{\mathrm{obs},0}}
 {\displaystyle\int_0^1 c\!\left(p_{A_0}(q)\right)\,dq}.
\end{equation}
Expected and realised curves each use their own multiplier. Making the expectation explicit,
the expected cost and captured gross alpha are
\begin{equation}\label{eq:costexpectation}
 \bar c(A)=K\int_0^1 c\!\left(p_{\mathrm{eff},A}(q)\right)\,dq,
\end{equation}
\begin{equation}\label{eq:alphaexpectation}
 R_g(A)=R_g(A_0)\int_0^1
 \frac{g(D_A(q),H_\alpha)}{g(D_0(q),H_\alpha)}\,dq.
\end{equation}
Here $g$ is equation~\eqref{eq:capture}, and a common half-life is used for the aggregate
model. These integrals are the notional-weighted expectations previously written as
$\mathbb E[\cdot]$. Integrating over uniform ranks is equivalent to integrating over the
participation distribution supplied by Section 3, with current and alternative horizons
paired at the same rank.

The expectations are evaluated numerically using $J$ equally weighted midpoint quantiles:
\begin{equation}\label{eq:quadrature}
 \begin{aligned}
 q_j&=\frac{j-\tfrac12}{J},\quad j=1,\ldots,J,\qquad
 K_J=\frac{\bar c_{\mathrm{obs},0}}
 {J^{-1}\sum_{j=1}^{J}c(p_{A_0}(q_j))},\\
 \bar c_J(A)&=\frac{K_J}{J}\sum_{j=1}^{J}c(p_{\mathrm{eff},A}(q_j)),\\
 R_{g,J}(A)&=\frac{R_g(A_0)}{J}\sum_{j=1}^{J}
 \frac{g(D_A(q_j),H_\alpha)}{g(D_0(q_j),H_\alpha)}.
 \end{aligned}
\end{equation}
The factor $1/J$ supplies the notional weight. There is no additional bucket-share multiplier.
Using the same numerical grid in the calibration denominator and cost average ensures
$\bar c_J(A_0)=\bar c_{\mathrm{obs},0}$ and $R_{g,J}(A_0)=R_g(A_0)$.
The continuous impact expectation requires $c\,d(A)>\gamma$ for the Burr tail; the finite-mean
participation constraint alone does not guarantee this for every fitted exponent or AUM.
Finite-grid runs should check sensitivity to $J$, particularly under tail thickening.

The resulting impact estimate is converted from cost per executed notional into annual
cost per unit of portfolio capital using one-way turnover:
\begin{equation}\label{eq:drag}
 C(A)=\frac{2T\bar c(A)}{10{,}000}.
\end{equation}
The conversion by $10{,}000$ changes basis points into a decimal return. Net alpha and IR are
\begin{equation}\label{eq:net}
 R_{\mathrm{net}}(A)=R_g(A)-C(A),\qquad
 IR_{\mathrm{net}}(A)=\frac{R_{\mathrm{net}}(A)}{TE}.
\end{equation}
The quantities $R_g$, $C$ and $TE$ are expressed as decimal annual returns, and $T$ as
decimal one-way turnover. Distributional migration therefore affects both components of
net return: impact in $C(A)$ and delay in $R_g(A)$. Tracking error and turnover remain fixed assumptions
in this specification.

For example, consider a notional rank whose participation increases from 8\% to 16\% ADV. With a 10\% daily
target and a non-binding horizon cap, its horizon increases from one to two days.
Effective participation becomes $16\%\times1/2=8\%$, so its modeled impact is unchanged,
while its alpha-capture ratio becomes $g(2,H_\alpha)/g(1,H_\alpha)<1$.
Migration can therefore lower net IR through delay even when scheduling offsets the impact
increase. This is an illustrative rank, not a claim about the entire portfolio.

A run report should show the complete capacity curve and explicitly selected thresholds rather
than presenting a single universal capacity number.

\section{Required run-report contents}
Every portfolio run report must include:
\begin{enumerate}
  \item run identity, timestamp, code version and input source;
  \item supplied strategy, liquidity, execution and alpha-decay inputs;
  \item fitted distribution and impact parameters with fit diagnostics;
  \item current-AUM reconciliation to observed costs and alpha;
  \item AUM scenario table and continuous capacity curves;
  \item bucket migration and execution-horizon diagnostics;
  \item selected capacity thresholds and exact crossing estimates;
  \item data-quality, extrapolation and model-risk warnings.
\end{enumerate}

\section{Validation requirements}
A valid run should confirm that notional shares sum to one, AUM and half-life are positive,
participation buckets are ordered and non-overlapping, expected and realised costs share a common
benchmark, and current-AUM impact reconciles to the supplied weighted cost. The report should flag
missing trade-level alpha weights, large extrapolations, horizons at $D_{\max}$, unstable tail fits
and sensitivity to user policy parameters.

\appendix
\section{User-supplied data inputs}
\small
\begin{longtable}{@{}p{0.23\textwidth}p{0.16\textwidth}p{0.18\textwidth}p{0.35\textwidth}@{}}
\caption{Required and optional user data}\label{tab:userinputs}\\
\toprule
Input & Requirement & Unit / format & Definition and validation\\
\midrule
\endfirsthead
\toprule
Input & Requirement & Unit / format & Definition and validation\\
\midrule
\endhead
Current AUM $A_0$ & Required & Currency bn & Average deployed AUM represented by the liquidity sample; must be positive.\\
Gross expected alpha $R_g$ & Required & Decimal p.a. & Expected relative return before modeled impact; document whether existing live slippage is embedded.\\
Tracking error $TE$ & Required & Decimal p.a. & Annualised ex-ante or realised relative volatility; must be positive and methodologically aligned with alpha.\\
One-way turnover $T$ & Required & Decimal p.a. & Purchases or sales divided by average AUM. Do not supply two-way turnover without changing the convention.\\
Average holdings $N$ & Diagnostic & Count & Average number of portfolio positions over the measurement period.\\
Bucket lower and upper bounds & Required for aggregate run & Fraction of ADV & Ordered, non-overlapping participation intervals covering the full positive range.\\
Executed-notional share & Required for aggregate run & Fraction summing to 1 & Annual executed-notional mass in each bucket; not trade-count frequency.\\
Representative participation & Required for impact fit & Fraction of ADV & Prefer the notional-weighted mean participation within each bucket.\\
Expected impact & Required & bp / executed notional & Expected implementation cost using a consistent arrival-price or other declared benchmark.\\
Realised impact & Recommended & bp / executed notional & Realised cost using the same benchmark as expected impact.\\
Parent-order notional or shares & Required for trade-level run & Currency / shares & Parent order before child-order slicing, net of cancellations according to a declared convention.\\
Security ADV & Required for trade-level run & Currency or shares/day & ADV measured consistently with order units, including declared lookback and lag.\\
Order timestamps & Recommended & Datetime & Decision, release, first fill and final fill times.\\
Actual execution horizon & Recommended & Trading days & Number of sessions over which the parent order was worked.\\
Forecast alpha & Recommended & bp or return & Expected contribution at the order decision timestamp.\\
Signal family & Recommended & Category & Used to assign or estimate signal-specific decay.\\
Observed alpha path & Preferred & Return by horizon & Forward alpha after 1, 2, 3, 5, 10 or more trading days for empirical half-life estimation.\\
Side, spread and volatility & Recommended & Category / bp / decimal & Used for richer asymmetric impact and execution-risk models.\\
\bottomrule
\end{longtable}

\section{User-set policy and scenario parameters}
\begin{longtable}{@{}p{0.25\textwidth}p{0.14\textwidth}p{0.17\textwidth}p{0.36\textwidth}@{}}
\caption{Parameters chosen or approved by the user}\label{tab:policy}\\
\toprule
Parameter & Example & Unit & Interpretation\\
\midrule
\endfirsthead
\toprule
Parameter & Example & Unit & Interpretation\\
\midrule
\endhead
Participation-scale elasticity $\eta$ & 0.85 & Elasticity & Response of the Burr scale to AUM. Set to 1 for proportional scaling; estimate from portfolio rescaling where possible.\\
Tail-thickening elasticity $\kappa$ & 0.00 & Elasticity & Structural change in upper-tail liquidity composition. Keep at zero unless supported by portfolio simulations or observations.\\
Target daily participation $\rho$ & 10\% & ADV/day & Execution-policy target used to determine horizon.\\
Maximum execution horizon $D_{\max}$ & 10 & Trading days & Maximum permitted horizon before daily participation exceeds the target.\\
Alpha half-life $H_\alpha$ & 10 & Trading days & Time for forecast alpha to halve. Prefer signal-specific empirical estimates.\\
AUM scenario grid & 1--100 & Currency bn & Range over which curves and thresholds are evaluated.\\
Net IR threshold & 0.75 & Ratio & Minimum acceptable net information ratio.\\
Alpha-retention threshold & 80\% & Percentage & Minimum fraction of current gross alpha remaining after delay and impact.\\
Impact-drag threshold & 100 & bp p.a. & Maximum acceptable annual return drag.\\
Alpha-capture threshold & 95\% & Percentage & Minimum gross alpha retained after execution delay, before impact.\\
\bottomrule
\end{longtable}

\section{Fitted parameters and outputs}
The user does not normally supply Burr parameters $(c,d,\lambda_0)$ or impact parameters
$(a,b,\gamma)$; the model estimates them. Every run report should disclose them along with the tail
index $cd$, bucket-fit residuals, impact-fit residuals, baseline calibration multiplier, threshold
crossings and warnings. Parameters should be versioned with their input sample and estimation date.
\normalsize

\begin{thebibliography}{9}\small
\bibitem{rak2007} R. Rak et al. ``Tests of scaling and universality of the distributions of trade
size and share volume.'' \emph{Physical Review E} 76, 046109 (2007).
\href{https://doi.org/10.1103/PhysRevE.76.046109}{doi:10.1103/PhysRevE.76.046109}.
\bibitem{farmer2011} J. D. Farmer et al. ``How efficiency shapes market impact.''
\emph{Quantitative Finance} 13(11), 1743--1758 (2013).
\href{https://arxiv.org/abs/1102.5457}{arXiv:1102.5457}.
\bibitem{almgren2001} R. Almgren and N. Chriss. ``Optimal Execution of Portfolio Transactions.''
\emph{Journal of Risk} 3, 5--39 (2001).
\href{https://doi.org/10.21314/JOR.2001.041}{doi:10.21314/JOR.2001.041}.
\bibitem{garleanu2013} N. G\^{a}rleanu and L. H. Pedersen. ``Dynamic Trading with Predictable
Returns and Transaction Costs.'' \emph{Journal of Finance} 68(6), 2309--2340 (2013).
\end{thebibliography}
\end{document}
